\section{Ascent to the group of order nine}\label{sec:ascent}

We now use the subgroup calculation to determine the fate of the
height-six detectors.  The argument takes place on finitely many pages
of the Euler-localized spectral sequence.  Its translation classes are
units on those pages; no integer-graded periodicity assertion about the
abutment is required.

\begin{theorem}\label{g:tail-exclusion}
For every $\jmath\geq4$, the class
$b_{\jmath}\in\Ext_{\calA_3}^{2,4\cdot3^{\jmath+1}}
(\F_3,\F_3)$ does not survive the classical Adams spectral sequence.
\end{theorem}

The proof will occupy this section.  The classical survivor $b_2$,
recalled in \cite[Corollary~1.2]{Amelotte}, will be
used to force an outgoing differential on an entire lattice of scalar
classes.  All the detectors with $\jmath\geq4$ belong to that lattice,
whereas the detector for $b_3$ does not.

\subsection{Coordinates, twists, and translation classes}

Retain $G=\langle g\rangle=C_9$, $H=\langle g^3\rangle=C_3$,
and the coefficient notation of Sections~\ref{sec:coefficients}
and~\ref{sec:subgroup}.  Put $Q=G/H$.  On $Q$-modules write
\[
 D=g-1,\qquad N=1+g+g^2.
\]
Thus $D$ here is a difference operator, not a deformation ring.  The
coefficient-field notation is $k=\F_{3^6}$.  Restriction and additive
transfer between the localized spectral sequences will be denoted
$\operatorname{res}$ and $\operatorname{tr}$.

Keep the initial $H$ exterior-degree coordinate $j\in\{0,1,2,3\}$.
The other coordinates are determined by
\begin{equation}\label{g:coordinates}
 s=2m+j,\qquad t_{\mathrm{int}}=4j-6n,\qquad
 w_0=m+2n+2j,\qquad r_0=w_0-j\pmod3.
\end{equation}
Here $s$ is cohomological filtration and $t_{\mathrm{int}}$ is internal
degree.  The parity of $s$ and the residue of $t_{\mathrm{int}}$ modulo
six determine $j$ uniquely among $0,1,2,3$.  Consequently the
coordinates distinguish all the types that occur.  We reserve
$\jmath$ for a Kervaire index.

Besides faithful Euler localization, include all powers of
$u'=u_{\lambda'}$, where $\lambda'=\lambda^3$ and $\lambda$ is the
faithful complex $G$-character used earlier.  The variable $u'$ is the
orientation symbol for the dimension-zero shift $2-\lambda'$.
Its restriction to $H$ is the trivial orientation.  We write $i\in\Z$ for
the total exponent of $u'$ in a homogeneous symbol.  Write
\[
 B=b^3,\qquad C=b^9,\qquad
 \Delta=\prod_{a=0}^8x_a.
\]
In particular $\operatorname{res}(\Delta)=q_0q_1q_2$ on cohomology.

\begin{lemma}\label{g:translations}
On every nonzero page the class
\begin{equation}\label{g:Delta-unit}
 \mathcal D=\Delta b^{-6}u'^2
\end{equation}
is a cycle and a unit.  Constants from $W(k)$ act by cycles.  There is
also a semilinear symmetry commuting with $G$, inducing Witt
Frobenius on $k$ and fixing $b$ and $u'$.
\end{lemma}

\begin{proof}
The multiplicative norm of the underlying class $x_0$ is an actual
representation-graded homotopy class.  In the real regular
representation of $G$, three of the nontrivial line pairs are
faithful, and the fourth is $\lambda'$.  A power map of complex lines
of degree prime to three gives an equivalence of their one-point
compactifications after three-localization.  Thus the three faithful
line pairs may all be identified with $\lambda$, with unit changes
of orientation.  The norm symbol is consequently
$\Delta u_\lambda^6u'^2$.  Division by the sixth power of the
faithful Euler symbol $a_\lambda=bu_\lambda$ gives
\eqref{g:Delta-unit}.  Its norm coefficient is invertible.  The
derivation rule shows that the inverse of a cycle unit is again a
cycle for as long as the page is nonzero.

All the Tate groups are annihilated by nine.  A Teichm\"uller
constant $c$ satisfies $c^{3^6}=c$; applying any differential gives
$dc=3^6c^{3^6-1}dc=0$.  Every Witt constant modulo nine is a sum of
a Teichm\"uller constant and three times one, proving the assertion
about constants.

Finally, Witt Frobenius on the maximal order is inner, implemented
by $S$, up to the chosen opposite-group convention.  A uniformizer
in the maximal subfield $\Q_3(\zeta_9)$ has the same valuation.
Multiplying $S$ by a unit therefore makes its conjugation agree
with conjugation by an element of that subfield.  The resulting
extended-stabilizer symmetry commutes with $G$.  It induces Witt
Frobenius on constants and preserves the ordinary cyclic and
representation-orientation symbols $b,u'$.
\end{proof}

When two differential coefficients occur in the same residue-field
target and their sources are fixed by this symmetry, their ratio,
if defined, belongs to $\F_3$.  We will use this elementary descent
observation repeatedly.

\subsection{The integral Mackey calculation}

Write $P_n=Kq^n$ for the scalar line on the initial $H$ page.  A
cyclic basis $a_i=g^ia_0$ of $V$ can be chosen with
\[
 \theta=a_0+a_1+a_2.
\]
Indeed $V$ has a semilinearly permuted basis, while $\theta$ is
invariant and primitive modulo the maximal ideal.  An invariant
primitive vector in that basis is a trace of a vector whose
coefficient is a unit, and its three translates form a basis by
Nakayama's lemma.  Set
\[
 \overline M=k[Q]/(1+g+g^2).
\]
Thus $V/K\theta=K\otimes\overline M$, with its semilinear action,
and $\bigwedge^2\overline M\cong k$.

\begin{lemma}\label{g:mackey}
Suppressing powers of $b$ and the orientation twist $u'$, the scalar
and exterior-degree-three groups at $G$ have a common description
$\widetilde P_n$.  Restriction fits into
\[
 0\longrightarrow \overline k_n\longrightarrow
 \widetilde P_n\xrightarrow{\operatorname{res}}P_n^Q
 \longrightarrow0,
 \qquad
 \overline k_n=
 \begin{cases}k,&3\mid n,\\0,&3\nmid n.
 \end{cases}
\]
Transfer embeds $(P_n)_Q$ in $\widetilde P_n$, with cokernel
$\overline k_n$.  Moreover
\begin{equation}\label{g:P-tate}
 \widehat H^d(Q;P_n)=\overline k_n
 \qquad(d\in\Z).
\end{equation}
For $3\mid n$, the exceptional class is represented by
$\Delta^{n/3}$; its transfer is the nonzero three-multiple in the
restriction kernel.  There is a class $c_G$ restricting to
$a_0a_1a_2$, and multiplication by $c_G$ identifies the scalar and
volume descriptions, including their exceptional kernels.

In exterior degrees one and two, restriction is injective onto all
invariants.  These groups are parametrized by transfer of
$P_na_0$ and $P_na_0a_1$, respectively.
\end{lemma}

\begin{proof}
We first work in the polynomial coefficient model of
Theorem~\ref{coef:model}, before localization and completion.  Group
monomials by their three degrees in the trace-zero blocks.  A
triplet whose entries are not all equal has a three-element orbit;
the corresponding module is induced from $H$.  The standard
restriction and transfer for an induced module give precisely its
invariant and coinvariant descriptions.

Consider an equal triplet $(d,d,d)$.  The module is tensor induction
of $\Sym^d L$, where
\[
 L=W(k)[H]/(N_H),\qquad N_H=1+\sigma+\sigma^2.
\]
To compute its stable type, use the exact multiplication sequence
obtained by quotienting the polynomial permutation representation
by its invariant linear sum.  In degree $d$ the polynomial
permutation module has Tate cohomology only from the monomial with
three equal exponents, and that occurs precisely when $3\mid d$.
The long exact Tate sequence therefore gives one $k$ in even
degrees, one $k$ in odd degrees, or zero, according as
$d\equiv0,1,2\pmod3$.

These Tate patterns determine the nonprojective lattice summands.
For completeness, a torsion-free $W(k)[H]$-lattice is a sum of
trivial, cyclotomic, and free lattices.  Map it by $(N_H,\sigma-1)$
into its invariant and trace-zero lattices.  These are free over
$W(k)$ and $W(k)[\zeta_3]$, respectively.  Modulo their maximal
ideals, the image glues two subspaces by an isomorphism.  Choosing
bases adapted to those subspaces splits the gluing into paired
one-dimensional graphs and unpaired summands.  The paired graph
is $W(k)[H]$, since
\[
 W(k)[H]=W(k)\mathbin{\times}_{k}W(k)[\zeta_3].
\]
This gives the asserted three lattice types.  Accordingly,
$\Sym^dL$ is stably $W(k)$, $L$, or zero in the three residue
classes of $d$.

Terms containing a free summand contribute nothing after tensor
induction.  For a nonconstant choice of factors this follows by
induction from $H$.  For repeated free factors, the diagonal
$H$-action on a tensor permutation basis is free; its $G$-action
is then free too, because every nontrivial subgroup of $G$
contains $H$.

Tensor induction of $L$ is stably the third suspension of the
trivial $G$-lattice.  To see this without losing the integral
extension, tensor the three short exact quotient sequences
defining $L$.  Every intermediate term is projective by the
preceding observation.  The cyclic permutation of the three
suspended factors has positive sign.  The two surviving ends thus
give a threefold stable suspension, whose odd Tate group is
$W(k)/9$.  Tensor induction of the trivial lattice gives the
corresponding even $W(k)/9$.  Restriction to $H$ is reduction
modulo three, and transfer is multiplication by three.  In
particular the exceptional transfer is nonzero; replacing these
groups by $k$ prematurely would lose the restriction kernel.

Here is the degree identification.  At $H$, degree
$3e+\epsilon$ in block $i$, with $\epsilon=0,1$, contributes
$q_i^ex_i^\epsilon$; the odd case is interpreted as a cyclic
one-cocycle value.  After allowing $\Delta^{-1}$, use
$x_i/q_i$ for the odd generator.  Then $j$ counts odd factors
and $n=\sum_i(e_i+\epsilon_i)$.  An equal triplet has $3\mid n$
and $j=0$ or $3$, giving exactly the exceptional terms above.
The odd equal-triplet generator supplies $c_G$.
The remaining triplets give all the degree-one and degree-two
invariants by transfer.  The same monomial calculation, now on
the three $q_i$, proves \eqref{g:P-tate}.

These calculations survive localization and completion.  After
inverting $\Delta$, the degree-zero coefficient ring is finite
over its Noetherian invariant subring, and every graded piece is
a finite module there.  Completion on the invariant side at the
image of the residue point is flat and gives the desired
completion, since the primes in the fibre constitute one orbit.
The exceptional terms are supported at that point: invariant
ratios act on them by their values at $x_i/x_0=1$ in $k$.
Flat completion of the periodic cyclic complexes therefore
preserves the computed groups and maps.  Finally, multiplication
by $c_G$ identifies their scalar and volume parts; on the
exceptional kernels this follows from the nonzero transfer
pairing, and on induced summands from restriction.
\end{proof}

\subsection{The pages through length nine}

\begin{lemma}\label{g:early-pages}
At fixed $r_0,n$, the $d_5$ complex in increasing exterior degree is
\begin{equation}\label{g:d5-complex}
 \widetilde P_n\xrightarrow{\ r_0\operatorname{res}\ }P_n
 \xrightarrow{\ \pm(r_0+1)D\ }P_n
 \xrightarrow{\ (r_0+2)\operatorname{tr}\ }\widetilde P_n.
\end{equation}
Consequently the $E_6$ page has the following rows:
\begin{equation}\label{g:E6}
\begin{array}{c|cccc}
 r_0&j=0&j=1&j=2&j=3\\\hline
 0&\widetilde P_n&P_n^Q&0&\overline k_n\\
 2&\overline k_n&P_n/P_n^Q&DP_n&\overline k_n\\
 1&\overline k_n&0&(P_n)_Q&\widetilde P_n.
\end{array}
\end{equation}
The orientation $u'$ has zero $d_5$ and $d_9$.
\end{lemma}

\begin{proof}
Theorem~\ref{h:differentials} gives the $H$ formula
$w_0\theta\wedge(-)$.  On successive exterior degrees in a row
of \eqref{g:d5-complex}, its coefficient is $r_0+j$.  Restriction
determines the first arrow.  Writing the cyclic basis modulo
$\theta$ identifies the middle arrow with $D$, up to an
invertible translate and sign.  Transfer determines the last
arrow by the projection formula.  Lemma~\ref{g:mackey} then
gives \eqref{g:E6}: in each of the three residue cases one
coefficient vanishes and the other two are units.  Restriction
of $u'$ is the trivial orientation; its possible $d_5$ target
has injective restriction.  At length nine it has no target.
\end{proof}

At $H$, the corresponding $E_6$ row has $P_n,\theta P_n$ in
row zero, the plane $P_n\otimes\overline M$ and its
$\theta$-multiple in row two, and its determinant and
$\theta$-multiple in row one.  Identify that determinant with
$k$ using $a_0a_1\bmod\theta$, which is fixed by $Q$.
The $d_9$ preserves $r_0$, decreases $j$ by one, and replaces
$n$ by $n-2$; its nonzero $H$ maps multiply a
$\theta$-multiple by $J$.

\begin{proposition}\label{g:E10-description}
Put
\[
 S_n=P_n/JP_{n+2},\qquad W_n=S_n\otimes\overline M.
\]
The main $H$ groups on $E_{10}$ are $S_n,W_n,S_n$ in
exterior degrees $0,1,2$, respectively, at $w_0\equiv0\pmod3$.
Define transfer modules
\begin{align}
 \mathcal T_n&=
 \begin{cases}(S_n)_Q,&n\equiv0,2\pmod3,\\
 NS_n,&n\equiv1\pmod3,
 \end{cases}\label{g:Tn}\\
 \mathcal W_n&=
 \begin{cases}(W_n)_Q,&n\equiv1,2\pmod3,\\
 NW_n,&n\equiv0\pmod3.
 \end{cases}\label{g:Wn}
\end{align}
At $G$ the main groups are:
\begin{enumerate}
\item in degree zero, a scalar group $\mathcal S_n$ with transfer
part $\mathcal T_n$ and one additional restriction generator
when $3\mid n$;
\item in degree one, the transfer group $\mathcal W_n$;
\item in degree two, the transfer group $\mathcal T_n$.
\end{enumerate}
On these transfer groups restriction is norm and transfer is the
displayed projection.  The additional groups are one $k$ each,
only when $3\mid n$: scalar restriction kernels at
$w_0\equiv1,2\pmod3$, and volume groups at every $w_0$.
The volume is of $3c_G$ kernel type when $w_0\equiv1$, and
of reduced $c_G$ type when $w_0\equiv0,2$.
\end{proposition}

\begin{proof}
In row zero of \eqref{g:E6}, the degree-one $d_9$ is injective
by restriction.  When its invariant source is a norm, the map
is also specified by transfer.  In row two both middle groups
inject under restriction, so the degree-two map is injective.
Its entire source is transferred from $\theta$-multiples,
represented by $DP_n$; the degree-one outgoing is therefore
zero by transfer.

For row one, choose $L_1$ with $NL_1=J$.  Such an element
exists by \eqref{g:P-tate}, since $J$ has index $-2$.
The top arrow sends the invariant restriction value $f$ to
$[L_1f]$ in $(P_{n-2})_Q$, up to a common nonzero scalar.
Restriction proves this except when the target index is
divisible by three.  In that case all the sources are
transfers, and
\[
 [L_1Nf']=[Jf']
\]
in coinvariants gives the same formula.  All remaining arrows
vanish by degree or by restriction.

Take the quotients row by row.  In degree zero, when
$n\equiv0$, every incoming invariant multiplier lifts by a
norm, so the transfer still embeds after passage to
$(S_n)_Q$.  When $n\equiv1$, restriction was injective on
the whole transfer target; quotienting its invariant
$J$-image leaves exactly $NS_n$.  The $n\equiv2$ case
has no additional Tate quotient and gives $(S_n)_Q$.
In degree one, restriction was injective before $d_9$;
the quotient is $(W_n)_Q$ divided by the reduction of
$\widehat H^{-1}(Q;P_n\otimes\overline M)$.  This removes
the norm kernel exactly for $n\equiv0$.  In degree two,
the additional quotient is by $L_1P_{n+2}^Q$; when
$n\equiv1$ it removes precisely the newly appearing norm
kernel.

For these assertions, the long exact Tate sequence of
multiplication by $J$ gives
\[
 \widehat H^d(Q;S_n)\cong\widehat H^d(Q;W_n)
 \cong
 \begin{cases}k,&n\equiv0,1\pmod3,\\0,&n\equiv2\pmod3,
 \end{cases}
\]
as statements about dimensions in every degree.  The plane
statement uses the stable suspension represented by
$\overline M$.  The same exact sequences identify the
norm kernels just used.  Finally, the endpoint kernels
and cokernels left in \eqref{g:E6} give the additional
scalar and volume lines stated in the proposition.
\end{proof}

\subsection{Restriction-zero corrections at lengths thirteen and seventeen}

Restriction does not determine a differential whose target is a norm
kernel.  We keep all such coefficients until naturality or
multiplication forces them.

\begin{lemma}\label{g:d13}
Put $T_0=b^5c_G$ on $E_{13}$.  After a nonzero normalization,
\[
 d_{13}B/B=\varepsilon T_0,\qquad
 d_{13}u'/u'=0,\qquad \varepsilon\ne0.
\]
There is a residue $\rho\in\{2,5,8\}\pmod9$ such that the extra
scalar families surviving $d_{13}$ occur in the two families
$w_0\equiv1\pmod3$ or $w_0\equiv\rho\pmod9$.
The extra volume families occur in the three residue families
\[
 w_0\equiv0\pmod3,\qquad 2\pmod9,
 \qquad \rho+2\pmod9.
\]
Scalar generators outside the transfer part survive precisely at
$w_0=9t$ and $3\mid n$.
\end{lemma}

\begin{proof}
At length thirteen only the exterior-degree change $0\to3$ can
be nonzero.  For the reverse change, the target index is $n-4$,
where there is no restriction kernel.  The main scalar-to-volume
map is necessarily zero on transfers, by its $H$ counterpart.
Write its two unit coefficients as
\[
 d_{13}B/B=\varepsilon T_0,\qquad
 d_{13}u'/u'=\varepsilon'T_0.
\]
For $w_0=3p$ and $3\mid n$, a scalar lift has symbol
\[
 \mathcal D^{n/3}B^pu'^{\,i-2n/3},
\]
and its logarithmic coefficient is
$p\varepsilon+(i-2n/3)\varepsilon'$.  If this vanished with
$p\not\equiv0\pmod3$, the restriction would survive to
$d_{17}$, where Theorem~\ref{h:differentials} gives a nonzero
primitive $h$ term.  But the available restriction image in the
$G$ target is $NW_{n-2}$, which is zero modulo the curve
parameter $Z$: the residual plane is a two-term Jordan block.
Naturality gives a contradiction.  Lemma~\ref{g:translations}
makes the ratio $\varepsilon'/\varepsilon$ Frobenius-fixed when
defined.  If $\varepsilon=0$, take $p=1$ and $i-2n/3=0$;
if $\varepsilon'\ne0$, choose $i-2n/3$ in $\F_3$ to cancel
the coefficient with $p=1$.  Both are impossible.  Thus
$\varepsilon\ne0$ and $\varepsilon'=0$.

On the two extra scalar families, the actual quotient Euler
symbol $3bu'$ gives $d_{13}(3b)=0$.  The product
$3b^2T_0$ is nonzero in its $3c_G$ target.  Hence we must
allow
\[
 d_{13}(3b^2)=\eta\varepsilon\,3b^2T_0,
 \qquad \eta\in\F_3.
\]
At $w_0\equiv1\pmod3$, multiplication by $T_0$ vanishes,
so the extra scalar remains.  At $w_0\equiv2\pmod3$,
the derivation rule for $B$ makes the scalar remain exactly
when $w_0\equiv\rho=2-3\eta\pmod9$.  Taking the
cokernels in the volume groups gives the three stated
families.  Finally the scalar lift coefficient is
$p\varepsilon$, so lifts outside transfers remain exactly
when $w_0=3p\equiv0\pmod9$.
\end{proof}

\begin{lemma}\label{g:d17}
On transfers, the length-seventeen chain anchored at a scalar
with $w_0=3p$ is
\begin{equation}\label{g:d17-exact}
 0\longrightarrow\mathcal T_n
 \xrightarrow{\quad h\quad}\mathcal W_{n-2}
 \xrightarrow{\quad h\wedge\quad}\mathcal T_{n-4}
 \longrightarrow0,
\end{equation}
with its arrows scaled by $p$ and $p+1$, respectively, up to
nonzero units.  The degree-two term has zero outgoing.

For a scalar lift at $w_0=9t$, $3\mid n$, the only possible
additional coefficient is
\begin{equation}\label{g:V-form}
 V(t,i-2n/3)=t\nu+(i-2n/3)\nu'
\end{equation}
into a one-dimensional restriction kernel in the plane.  The
ratio of $\nu,\nu'$ is Frobenius-fixed whenever defined.
\end{lemma}

\begin{proof}
At $H$, the invariant line generated by $h$ and its wedge
quotient give
\[
 0\longrightarrow S_n\longrightarrow W_{n-2}
 \longrightarrow S_{n-4}\longrightarrow0.
\]
Apply coinvariants and their long exact Tate sequence.
When $n\equiv0$, the last Tate group is zero, giving
\eqref{g:d17-exact} directly.  When $n\equiv1$, the
middle Tate group is zero and exactness requires dividing
the first coinvariant group by its norm kernel, exactly as
in \eqref{g:Tn}.  When $n\equiv2$, the first Tate group
is zero, and the corresponding norm kernels in the middle
and last coinvariants are both divided out.  These are
all the cases.  Equivalently the obstruction is
\[
 \ker(N:M_Q\longrightarrow M)
 =\widehat H^{-1}(Q;M),
\]
which has the dimensions already calculated.  Transferring
the $H$ differentials now gives the two coefficients in
\eqref{g:d17-exact}; the shift between the two source
coordinates changes $p-1$ in the $H$ formula to $p+1$
in this anchored chain.  The top outgoing is zero by
transfer.

On scalar unit lifts the restriction of a possible extra
map is zero.  Its target is the single norm kernel in the
first unscaled image of \eqref{g:d17-exact}, occurring
when $p\equiv0$ and $3\mid n$.  Derivation on $C,u'$
and translation by $\mathcal D$ give exactly
\eqref{g:V-form}.  There is no other kernel target.
The Frobenius assertion follows from
Lemma~\ref{g:translations}.
\end{proof}

For later use we record the full $E_{18}$ description.  The
extra scalar and volume families of Lemma~\ref{g:d13} are
unchanged.  The main terms are
\begin{equation}\label{g:E18}
\begin{array}{c|c|p{9.1cm}}
 j&w_0\bmod9&\text{group}\\\hline
 0&0&$\mathcal T_n$, together with the extra lift when
 $3\mid n$ and $V(t,i-2n/3)=0$;\\
 1&6&$\mathcal T_{n+2}$, identified by $h$; divide out its
 norm kernel if $3\mid n+2$ and
 $V(t,i-2(n+2)/3)\ne0$;\\
 1&3&$\mathcal T_{n-2}$, identified by $h\wedge$;\\
 2&0&$\mathcal T_n$.
\end{array}
\end{equation}
Here and below residues are represented by $0,\ldots,8$, and
$t=(w_0-(w_0\bmod9))/9$.  This table follows from
\eqref{g:d17-exact}: for $p\equiv0$ its scalar and first
image remain; for $p\equiv1$ the complex is acyclic;
for $p\equiv2$ its quotient and square remain.  The
correction \eqref{g:V-form} removes exactly the stated
norm kernel.  The maps called restriction and transfer
remain the norm and projection maps in these descriptions.

\subsection{Reduction to a finite array}

The main length-thirty-three contractions remove the curve parameter
from the surviving transfer modules. Some restriction-zero arrows of
length at most thirty-three remain undetermined. The finite array below
therefore retains their possible endpoints, subject to the constraints
derived later.

The degree changes needed below follow immediately from
\eqref{g:coordinates}.  They are displayed together so that
all exclusions can be checked without an implicit chart:
\begin{equation}\label{g:raw-lengths}
\begin{array}{c|c|c|c}
 \Delta j&\text{length}&\Delta w_0&\Delta n\\\hline
 1&12l+5&2l+4&-2l\\
 -1&12l+9&2l-1&-2l-2\\
 3&12l+1&2l+9&2-2l\\
 -3&12l+1&2l-8&-2-2l.
\end{array}
\end{equation}
They exhaust the possibilities because a differential changes
the parity of exterior degree, and $0\le j\le3$.

\begin{lemma}\label{g:to33}
After the main length-thirty-three maps, the only potential
groups are those in Table~\ref{g:finite-table} below, subject
to the qualifications following the table.
\end{lemma}

\begin{proof}
First consider the possible intervening lengths in
\eqref{g:E18}.  At length twenty-one, the degree-one sources
are transfers and their $H$ differential is zero; the
other arrows are excluded by \eqref{g:raw-lengths}.
At twenty-five the possible volume-to-scalar arrows have
source residues
\[
 2,\quad\rho+4,\quad\rho+2
\]
and target residues $7,\rho,\rho-2$, respectively.
The last target occurs only if $\rho=2$.  Any
degree-compatible scalar-to-volume alternative has
transfer source and is zero.  At twenty-nine only the
extra scalar at residue $4$ can hit a restriction kernel
in degree one at residue $3$: the other candidate has
target with injective restriction and is zero.  At
thirty-three the $H$ contractions give the main transfer
maps, and there may additionally be a volume arrow from
residue $6$ to the restriction kernel in degree two at
residue $0$.

It remains to compute the kernels and cokernels of the
main contractions.  Set $L=h^3$, an invariant scalar
section of index $-8$ and valuation one in the curve
ring $S_0=k[[Z]]$.  The $Q$-action on this disc is
nontrivial: otherwise its action on the scalar lines
would also be trivial, since a scalar of order three
in the unit group of the characteristic-three disc
is one.  This contradicts their calculated Tate
dimensions.  Thus its invariant fraction field has
degree three, and the complete disc extension is
totally ramified.  For example the orbit product of a
parameter generates an invariant disc over which the
original disc has rank three.  Valuations of invariant
functions are consequently multiples of three.

The section $L$ shifts index by $-8\equiv1\pmod3$,
while $\Delta$ shifts it by three and is an invariant
unit.  These sections realize minimal invariant
valuations $0,1,2$ at indices $n\equiv0,1,2$,
respectively.  Any ratio of two invariant sections of
one index is an invariant function, so there can be no
smaller nonnegative valuation in that index.  The Tate
dimensions of Proposition~\ref{g:E10-description}
then show that the corresponding norm images begin at
valuations $3,4,2$.

It follows that multiplication by $L$ on
$\mathcal T_n\to\mathcal T_{n-8}$ has the three
effects
\begin{equation}\label{g:L-cases}
\begin{array}{c|l}
 n\bmod3&\text{effect}\\\hline
 0&\text{surjective, with the whole norm kernel as kernel;}\\
 1&\text{injective, with cokernel }k;\\
 2&\text{injective, with cokernel the target norm kernel }k.
\end{array}
\end{equation}
Indeed multiplication shifts each free invariant lattice
by one valuation, and comparison with the minima
$3,4,2$ gives the stated free cokernels; the remaining
kernel or cokernel is the explicit norm kernel in
\eqref{g:Tn}.  Dividing a norm kernel out at length
seventeen simply removes that kernel or cokernel from
this list.

More precisely, the degree-one line at index $n$ is
identified by $h$ with $\mathcal T_{n+2}$, and its
contraction is
\[
 L:\mathcal T_{n+2}\longrightarrow\mathcal T_{n-6}.
\]
The degree-two group at index $n$ contracts to the
degree-one quotient identified by $h\wedge$, by
\[
 L:\mathcal T_n\longrightarrow\mathcal T_{n-8}.
\]
These formulas are transfers of the nonzero $H$
length-thirty-three maps.  Applying
\eqref{g:L-cases} to each index gives exactly the
finite table.  The extra volume arrow identified above
is the only additional incoming on the residual
degree-two kernel.
\end{proof}

We use sans-serif names for entries of the finite array, to
distinguish them from the Laurent classes $A,B,C$.  Put
$n_0=n\bmod3$ and $w_*=w_0\bmod9$.
\begin{table}[htbp]
\centering
\begin{tabular}{c|c|c|c}
\toprule
label&$j$&$n_0$&$w_*$\\\midrule
$\mathsf A,\mathsf R$&0&0,2&0\\
$\mathsf p,\mathsf q,\mathsf r$&0&0&1,4,7\\
$\mathsf v$&0&0&$\rho$\\
$\mathsf x,\mathsf X$&1&1,2&3\\
$\mathsf y$&1&1&6\\
$\mathsf M$&2&0&0\\
$\mathsf a,\mathsf b,\mathsf c$&3&0&0,3,6\\
$\mathsf d,\mathsf f$&3&0&$2,\rho+2$\\\bottomrule
\end{tabular}
\caption{Potential finite groups after the main $d_{33}$.
Paired lists are read entry by entry.  The residue $\rho+2$
is reduced modulo nine.}\label{g:finite-table}
\end{table}

Here $\mathsf A$ contains its transfer kernel unless it was
hit at twenty-five, and independently contains a restriction
lift when $V=0$.  We do not split that extension.  The groups
$\mathsf R,\mathsf x,\mathsf X$ are one $k$ each,
are transfers, and have zero restriction; $\mathsf X$ is
absent if hit at twenty-nine.  The group $\mathsf y$ is
the residual line kernel just when the indicated $V=0$;
$\mathsf M$ is the residual square kernel unless hit at
thirty-three.  The remaining entries are the extra scalar
and volume lines, with the endpoints of any nonzero
twenty-five, twenty-nine, or thirty-three arrows removed.
Every individual entry except the possible extension
$\mathsf A$ is one-dimensional over $k$.

Transfers of the $H$ constants generate the transfer parts of
$\mathsf A,\mathsf R,\mathsf x,\mathsf X$ whenever
present.  This follows from the cokernels in
\eqref{g:L-cases}, where transfer still surjects from
constants modulo $Z$.  Restriction is zero because the
norm images have positive valuation.  The lift in
$\mathsf A$, when present, restricts to a nonzero
constant.  For $\mathsf y,\mathsf M$, the relevant
norm kernel before $d_{33}$ is represented by a constant
in $(S_e)_Q$ with $e\equiv0\pmod3$, using $\Delta$
to shift indices.  Reduction at $Z=0$ shows that this
coinvariant constant is nonzero.

We shall also use products on the finite array.  The extra
scalar classes are shifts of three, and the volume classes
are shifts of the same odd class $c_G$.  Therefore products
of two extra scalars or two volumes are zero.  Their mixed
products retain the nonzero raw three-multiple pairing
whenever the factors and the target $3c_G$ line
$\mathsf f$ are still present.  This follows already on
$E_2$ from Lemma~\ref{g:mackey} and continues to hold
on each page on which those classes remain.

Index all copies of a label by
\begin{equation}\label{g:copy-index}
 t=(w_0-w_*)/9,\qquad u=i-2\lfloor n/3\rfloor.
\end{equation}
Together with $(j,n_0,w_*)$, these indices determine its
position modulo the permanent translation $\mathcal D$.

\subsection{The normalizer character and an exhaustive degree rule}

\begin{lemma}\label{g:character}
On every residue-field entry of Table~\ref{g:finite-table},
including both subquotients of $\mathsf A$, a normalizing
symmetry acts by
\begin{equation}\label{g:character-formula}
 d_0^{\,2m+3n+\phi},\qquad d_0^2=-1,
\end{equation}
up to the same orientation-twist factor at fixed $i$, where
\[
 \phi=0\quad\text{on scalar labels and }\mathsf x,\mathsf X,
 \qquad
 \phi=2\quad\text{on all other labels}.
\]
Consequently a differential of length $R$ can join a source
and target only if
\begin{equation}\label{g:character-test}
 R-\Delta j+3\Delta n+\phi_{\mathrm{target}}
 -\phi_{\mathrm{source}}\equiv0\pmod4.
\end{equation}
\end{lemma}

\begin{proof}
By Skolem--Noether choose a stabilizer unit $\delta$
conjugating $g$ to $g^{-1}$; a conjugator is made a unit
by adjusting its valuation in the ramified centralizing
field.  If $d_0$ is its residual Lie scalar, conjugation
scales the first strict coefficient of $g$ by $d_0^{-2}$,
whereas inversion negates that coefficient.  Thus
$d_0^2=-1$.

A cyclic $b$-shift changes sign and contributes exponent
two.  A scalar $q^n$ has residual Lie exponent $3n$.
Let $\overline V=V/\mathfrak nV$ and let $N_0=g-1$ on
$\overline V$.  The quotient $\overline V/N_0\overline V$
is represented by $x_0/q_0$.  Its coefficient weight $-2$
and the cyclic one-cocycle inversion sign give total
exponent zero.  Since
\[
 \delta N_0\delta^{-1}=-N_0+N_0^2,
\]
$\delta$ preserves the canonical Jordan flag and its
successive diagonal characters, starting at this quotient,
are $1,-1,1$.  The last line is
$k\theta=N_0^2\overline V$.  Thus on
$\overline W=\overline V/k\theta$ the quotient character
is $1$ and the invariant-line character is $-1$.
The primitive line generated by $q^2h$ reduces to this
unique invariant line, so its extra exponent is two.
The plane and volume determinants likewise have
character $-1=d_0^2$.

Restriction and transfer preserve these characters on
the finite quotients.  In particular the constants
representing the norm kernels $\mathsf y,\mathsf M$
have the indicated determinant character.  Their constant
representatives inject into the residue at $Z=0$, so
higher-order semilinear terms do not alter these characters.  On
$\mathsf A$ only its two subquotient characters are
needed; no splitting of their extension is asserted.
Conjugation identifies the representation spheres
with the same orientation factor on source and target
at fixed $i$.  Finally $2\Delta m=R-\Delta j$ by
\eqref{g:coordinates}, giving
\eqref{g:character-test}.
\end{proof}

Here is a convenient complete arithmetic procedure for all
the arrow tables that follow.  For $d_1=\Delta j$, define
\begin{equation}\label{g:arrow-rule}
\begin{array}{c|rrrr}
 d_1&1&-1&3&-3\\\hline
 c_*&5&9&1&1\\
 h_*&4&-1&9&-8\\
 k_*&0&-2&2&-2.
\end{array}
\end{equation}
Then $R=12l+c_*$, $\Delta w_0=2l+h_*$,
$\Delta n=-2l+k_*$.  For any ordered pair of labels,
the necessary and sufficient arithmetic tests are
\begin{align}
 l&\equiv5(w_{*,\mathrm{target}}-w_{*,\mathrm{source}}-h_*)
       \pmod9,\label{g:l-congruence}\\
 n_{0,\mathrm{target}}&\equiv n_{0,\mathrm{source}}-2l+k_*
       \pmod3,\label{g:n-congruence}\\
 2l+c_*-d_1+3k_*+\phi_{\mathrm{target}}
       -\phi_{\mathrm{source}}&\equiv0\pmod4.
       \label{g:four-congruence}
\end{align}
If they have a solution, they determine $l$ modulo eighteen
and therefore $R$ modulo $216$.  They determine the exact
copy shifts by
\begin{equation}\label{g:copy-shifts}
 \Delta t=\frac{w_{*,\mathrm{source}}+\Delta w_0
                    -w_{*,\mathrm{target}}}{9},
 \qquad
 \Delta u=-2\left\lfloor
       \frac{n_{0,\mathrm{source}}+\Delta n}{3}
                    \right\rfloor.
\end{equation}
Adding $216$ to a length adds $(4,24)$ to these shifts.
In all tables an arrow is considered only when both of its
entries are still present.

\subsection{Forcing the two length-fifty-three coefficients}

\begin{proposition}\label{g:rho-and-e53}
The residue in Lemma~\ref{g:d13} is $\rho=8$.
The coefficient $e_{53}$ of the $H$ length-fifty-three
differential is nonzero.  Moreover, in \eqref{g:V-form}
one has $\nu\ne0$ and $\nu'=0$.
\end{proposition}

\begin{proof}
Apply \eqref{g:l-congruence}--\eqref{g:four-congruence}
to the scalar entry $\mathsf A$.  Its only incoming
possibilities, modulo $216$, are from $\mathsf f$ at
length $25,169,97$ when $\rho=2,5,8$, respectively.
Its outgoing possibilities are at $13,17,53$, to
$\mathsf d,\mathsf y,\mathsf x$, respectively.
For clarity the complete incoming list within the Euler
bound is
\begin{equation}\label{g:A-incoming}
\begin{array}{c|l}
 \rho&\text{incoming lengths on }\mathsf A\text{ through }469\\\hline
 2&25,241,457\\
 5&169,385\\
 8&97,313.
\end{array}
\end{equation}
At length twenty-five this is a statement about the
residual transfer kernel after the main length-thirty-three
quotient, and may equivalently be read on the early page.
If $\rho=2$, that arrow is zero.  Indeed
$\mathsf f=\mathsf v\mathsf d$ up to shifts,
$\mathsf v$ has no outgoing there, and the possible
$\mathsf d$ differential lands in $\mathsf r$,
whose product with $\mathsf v$ is zero.

Now take an $\mathsf A$ copy with $e_{53}t=0$.
Its transfer-kernel generator is the transfer of an
$H$ scalar that survives $d_{53}$ and is a
$d_{105}$ target, by Proposition~\ref{h:unit-death}.
The corresponding $H$ source has exterior degree one
and index divisible by three.  Its transfer is zero
in Table~\ref{g:finite-table}.  Naturality of $d_{105}$
therefore forces the transferred scalar to vanish
\emph{before} the $G$ page $E_{105}$: on that page its
transfer would have to equal $d_{105}(0)=0$.
It cannot have an outgoing earlier, since it is
transferred from an $H$ cycle.  An incoming is thus
necessary before length $105$.  The $\rho=5$ row of
\eqref{g:A-incoming} has none, and the $\rho=2$
candidate is zero.  Hence $\rho=8$ and the required
incoming has length $97$.

If $e_{53}=0$, this reasoning applies to every copy of
the transfer kernel, and every $\mathsf f$ is consumed
at length $97$.  That differential cannot hit a
restriction lift of the unit, since $H$ restriction of
the target is nonzero.  No source then remains for the
only later possible unit incoming, at $313$.
This contradicts the bound of Proposition~\ref{det:euler-bound}.
Therefore $e_{53}\ne0$.

An $\mathsf A$ restriction lift with $V(t,u)=0$
survives to length fifty-three.  If $t\not\equiv0\pmod3$,
its restriction has nonzero $H$ differential, while the
only $G$ target is $\mathsf x$, whose restriction is
zero.  Thus $V(t,u)=0$ implies $t=0$ modulo three.
The Frobenius ratio condition in Lemma~\ref{g:d17}
makes this a linear equation over $\F_3$.  Its kernel
can be contained in $t=0$ only if $\nu\ne0$ and
$\nu'=0$.
\end{proof}

It follows that the entry $\mathsf y$ occurs exactly when
$t\equiv0\pmod3$.  Through $E_{53}$, multiplication by
$C^3$ and $u'$ gives unit symmetries, so all previous
removals are uniform for a fixed residue $t\bmod3$.
Transfer of the $H$ differential at length fifty-three
has two more consequences:
\begin{enumerate}
\item the transfer kernel in $\mathsf A$ maps to
$\mathsf x$ when $t\not\equiv0\pmod3$;
\item $\mathsf X$ can remain only at $t\equiv1\pmod3$.
Indeed the $H$ scalars with index $n\equiv1$ have
zero transfer, while their nonzero length-fifty-three
outgoing would otherwise transfer to that entry.
\end{enumerate}
The only earlier possible removal of $\mathsf X$ was
the length-twenty-nine arrow from $\mathsf q$.
Accordingly $\mathsf q$ and $\mathsf X$ have a
common survivor flag $Q_0\in\{0,1\}$, at residues
$t=0$ and $t=1$, respectively.

On scalar lifts with $t=3T$, the remaining length-fifty-three
coefficient into $\mathsf x$ has the form
\begin{equation}\label{g:omega-form}
 T\omega+u\omega',
\end{equation}
by derivation on $C^3,u'$ and translation by
$\mathcal D$.  The ratio is again Frobenius-fixed.
The nonvanishing of $\omega$ requires a norm argument;
restriction alone cannot prove it.

\begin{lemma}\label{g:finite-comparison}
On page $E_R$, ordinary-to-Tate comparison is surjective
in filtration $s\ge0$ and injective in filtration
$s\ge R-1$, in every representation twist used here.
If a specified ordinary filtration-zero symbol is a Tate
cycle through length $R-1$, then it extends as an
ordinary class over the $(R-1)$-skeleton.
\end{lemma}

\begin{proof}
At $E_2$ the comparison is the usual map from group
cohomology to Tate cohomology: it is surjective in
filtration zero and an isomorphism in positive
filtration.  Suppose the stated ranges hold on $E_R$.
Lift a Tate cycle of filtration $s\ge0$ to an ordinary
class.  Its differential maps to zero, and the target
has filtration $s+R\ge R$, where comparison is
injective.  Thus the lift is an ordinary cycle.
For injectivity on the next page, an ordinary cycle in
filtration $s\ge R$ that is a Tate boundary has a
Tate boundary source in filtration $s-R\ge0$.
Lift that source by surjectivity and use injectivity
on $E_R$ at filtration $s$ to identify its ordinary
boundary with the given cycle.  This proves the
induction.

Apply it to a specified filtration-zero symbol.  At
each earlier length its ordinary differential maps
to zero, and comparison is injective at that target;
thus that very symbol survives.  The finite tower
description of cycles supplies its lift over the
stated skeleton.  No identification of an infinite
Tate abutment enters this argument.
\end{proof}

\begin{lemma}\label{g:omega-norm}
In \eqref{g:omega-form}, one has $\omega\ne0$.
\end{lemma}

\begin{proof}
Consider at $H$ the ordinary orientation symbol
\[
 Y=u_{\lambda|H}^{-9}.
\]
On Tate pages it is $C$ times an inverse faithful
Euler power.  By Proposition~\ref{g:rho-and-e53} its
first nonzero differential has length fifty-three.
Lemma~\ref{g:finite-comparison} supplies an actual
$H$-equivariant representative on the $52$-skeleton
of the chosen free $G$-model for $EG$.
This is a lift of the specified filtration-zero symbol:
at each preceding length its ordinary differential
vanishes because comparison is injective at that
outgoing target.  We use this same skeleton for the
pointwise norm, restriction, and transfer.

Take its pointwise multiplicative norm on that same
skeleton.  Explicitly, tensor-induce its map, precompose
with the diagonal of the skeleton, and use indexed
multiplication into $E$.  Include its virtual sphere
factor before taking the norm.  This gives induced
degree
\[
 9\operatorname{Ind}_H^G(\lambda|H)
       -18(1+\lambda'),
\]
whose orientation symbol is
\[
 u_\lambda^{-27}u'^{18}
\]
up to a constant unit.  Here the three faithful lifts
of $\lambda|H$ have been identified by the prime-to-three
power equivalences used in Lemma~\ref{g:translations}.

We justify the first obstruction of this norm
directly.  On the boundary of a free $53$-cell the
partial maps define a product over a $52$-sphere.
Stably, the disjoint-basepoint sphere splits into a
constant summand and a reduced sphere.  The reduced
diagonal into two or more reduced factors maps
$S^{52}$ into $S^{104}$ or a higher-dimensional sphere,
and is zero.  Hence only terms with one reduced
factor contribute to the first obstruction.  The
other factors are their filtration-zero point values.
Sum these single-factor terms over the three
$H$-cells: the result is precisely the additive
transfer cochain of the one-factor obstruction,
multiplied by the other two translated symbols.
All three partial factors extend over the preceding
skeleton, so this formula is valid modulo exactly
the previous obstruction indeterminacies.  It is
therefore a formula on $E_{53}$:
\begin{equation}\label{g:norm-obstruction}
 d_{53}\operatorname{Nm}_H^G(Y)
 =\operatorname{tr}_H^G\bigl(d_{53}Y\cdot gY\cdot g^2Y\bigr).
\end{equation}
This local obstruction calculation is the filtration-zero
instance of the general norm-transfer differential
construction in \cite[Theorem~4.6 and
Proposition~5.4]{DHLLSWX}.  That construction is an
antecedent for the argument; its prime-two periodicity
applications are not inputs here.

The translates must in general be retained in
\eqref{g:norm-obstruction}.  In this application $Y$
has coefficient one and its $H$ representation extends
to $G$.  Its cohomology symbol is thus fixed by $Q$,
and its successive surviving page classes are fixed
as well.  Consequently the right side is
$\operatorname{tr}(Y^2d_{53}Y)$ here.

Multiply by the necessary faithful Euler power to
remove the $u_\lambda$ shift.  The source becomes
$C^3u'^{18}$, and the target becomes
\[
 \operatorname{tr}\bigl(C^2d_{53}^H(C)\bigr).
\]
Before reduction of indices, this target has
$n=-8$ and $w_0=39$.  It is the $\mathsf x$ entry;
that entry has no possible removal before length
fifty-three and receives the $H$ transfer
isomorphically on this page.  Explicitly its exterior
degree is one, its twist is $i=18$, and
\eqref{g:copy-index} gives $(t,u)=(4,24)$.
The complete arrow rule gives no incoming to
$\mathsf x$ before $53$; the Mackey calculation and
the main $d_{33}$ quotient identify the transferred
$H$ constant with its one-dimensional generator.
This assertion is on $E_{53}$, before taking the
cohomology of the differential now being computed.
The target is therefore
nonzero.  On the source $C^3u'^{18}$, formula
\eqref{g:omega-form} reads
$\omega+18\omega'=\omega$.  This proves the lemma.
\end{proof}

Put $\kappa=\omega'/\omega\in\F_3$ and define
\begin{equation}\label{g:lattice}
 \mathcal L=\{(t,u)\in\Z^2:t+3\kappa u\equiv0\pmod9\}.
\end{equation}
The residual $\mathsf A$ restriction lifts and the
$\mathsf x$ support after length fifty-three are,
respectively,
\begin{equation}\label{g:A-x-support}
 \mathcal L,\qquad(1,6)+\mathcal L.
\end{equation}
From now on a subscript on a finite label denotes $t\bmod3$, not a
power.  A statement about such a residue concerns every copy unless
a particular copy is specified.
The shift $(1,6)$ follows from
\eqref{g:copy-shifts}.  For now all transfer kernels of
$\mathsf A$ at $t\equiv0\pmod3$ remain independently
of those lifts.  The length-$105$ argument in
Proposition~\ref{g:rho-and-e53} forces every
$\mathsf f$ at $t\equiv2\pmod3$ to survive until
its nonzero length-$97$ map onto one of these kernels;
the shift there is $(1,12)$.  We call these the
\emph{required} $\mathsf f_2$ copies.

\subsection{Products and the first finite pairings}

The complete possible arrows from lengths $25$ through
$97$, apart from the main $d_{33}$ already taken, are
\begin{equation}\label{g:early-finite-arrows}
\begin{array}{c|l}
25&\mathsf b\longrightarrow\mathsf v,\quad
    \mathsf d\longrightarrow\mathsf r\\
29&\mathsf q\longrightarrow\mathsf X\\
33&\mathsf c\longrightarrow\mathsf M\\
49&\mathsf a\longrightarrow\mathsf R\\
53&\mathsf A\longrightarrow\mathsf x\\
57&\mathsf x\longrightarrow\mathsf p\\
61&\mathsf R\longrightarrow\mathsf f\\
85&\mathsf p\longrightarrow\mathsf c,\quad
    \mathsf q\longrightarrow\mathsf a,\quad
    \mathsf r\longrightarrow\mathsf b\\
93&\mathsf y\longrightarrow\mathsf p\\
97&\mathsf a\longrightarrow\mathsf v,\quad
    \mathsf d\longrightarrow\mathsf p,\quad
    \mathsf f\longrightarrow\mathsf A.
\end{array}
\end{equation}
This table is obtained by substituting the fifteen labels
of Table~\ref{g:finite-table}, with $\rho=8$, in
\eqref{g:l-congruence}--\eqref{g:four-congruence}; thus
it lists possibilities, not asserted nonzero maps.
The length-fifty-seven arrow is zero by transfer: a
remaining $\mathsf x$ is transferred from the $H$ slot
that survives fifty-three, since its $t$ residue is one.

The mixed products required below are
\begin{equation}\label{g:finite-products}
 \mathsf p\mathsf a=\mathsf f,\qquad
 \mathsf q\mathsf c=\mathsf f,\qquad
 \mathsf v\mathsf d=\mathsf f,\qquad
 \mathsf r\mathsf b=\mathsf f,
\end{equation}
up to nonzero units and with aligned degree shifts.
The $t$ indices add in the first pairing.  In each of
the other three pairings they add with an additional
one, since the sum of their $w_*$ representatives
exceeds the target representative by nine.  Products
of two extra scalars, or of two volumes, are zero.

\begin{lemma}\label{g:d25-nonzero}
Both length-twenty-five arrows in
\eqref{g:early-finite-arrows} are nonzero on every
copy.  Thus $\mathsf b,\mathsf v,\mathsf d,\mathsf r$
are all removed at that length.
\end{lemma}

\begin{proof}
Choose arbitrary copies $z$ of $\mathsf b$ and $z'$
of $\mathsf d$.  Their product is zero, so
\[
 0=d_{25}(zz')=d_{25}(z)z'-z\,d_{25}(z').
\]
Both arrows have shift $(-1,4)$ by
\eqref{g:copy-shifts}.  The two terms, when nonzero,
belong to the same $\mathsf f$ copy, at
$(t_z+t_{z'},u_z+u_{z'}+4)$, and the products in
\eqref{g:finite-products} are nonzero.  Thus one
coefficient vanishes if and only if the other does.
Because the two chosen copies were arbitrary, either
both arrows vanish everywhere or both are nonzero
everywhere.  This proves uniformity even before any
smaller unit lattice is used.

Suppose they vanish.  Then all $\mathsf v,\mathsf d$
copies persist to page $97$: they have no intervening
incoming, $\mathsf v$ has no outgoing, and the only
remaining possible outgoing of $\mathsf d$ is the
length-$97$ map to $\mathsf p$.  Choose their product
in a required $\mathsf f_2$ copy.  On that page,
\[
 d_{97}(\mathsf v\mathsf d)
   =\mathsf v\,d_{97}(\mathsf d)=0,
\]
because its possible value is the product of two
extra scalars.  This contradicts the required
nonzero differential on $\mathsf f_2$.  Therefore
both arrows are everywhere nonzero; their source
and target groups are one-dimensional, so all four
labels disappear.
\end{proof}

Two particular $\mathsf p$ copies have no outgoing while
they remain.  The first is the quotient Euler symbol
$3bu'$ at
\begin{equation}\label{g:Euler-copy}
 (t,u)=(0,1).
\end{equation}
The second is the detector for the known surviving
$b_2$, at
\begin{equation}\label{g:b2-copy}
 (t,u)=(-4,12).
\end{equation}
We denote them $p^E$ and $p^\dagger$, respectively;
their $t$ residues are zero and two.

Here is the detector coordinate calculation, including its
cycle assertion.  By Theorem~\ref{det:actual-detector},
a survivor $b_{\jmath}$, $\jmath\ge2$, has nonzero
ordinary $G$ detection in filtration two.  At $H$ its
scalar coordinate is $m=1$, $n=-2\cdot3^{\jmath}$,
so every nonzero main scalar restriction has a nonzero $d_5$:
coefficients in $K$ are cycles and multiplication by the
primitive vector $\theta$ is injective, even for arbitrary
completed coefficients.  Since a sphere lift is an ordinary cycle, its
restriction must therefore vanish.  The $G$ image is
the generator of the extra scalar restriction kernel,
with $w_0=1-4\cdot3^{\jmath}\equiv1\pmod9$.
There is no incoming in ordinary filtration two; its
Tate image has zero outgoing on every page where it
has not yet been hit.  Formula~\eqref{g:copy-index}
now gives
\begin{equation}\label{g:detector-coordinates}
 (t,u)=
 \bigl(-4\cdot3^{\jmath-2},\,4\cdot3^{\jmath-1}\bigr).
\end{equation}
This explains \eqref{g:b2-copy} and the two cycle
assertions.  It does not assert that a sphere detector
cannot be an incoming Tate boundary.

\begin{lemma}\label{g:a-removal}
The entries $\mathsf a_0,\mathsf a_2$ and their
paired $\mathsf R_0,\mathsf R_2$ are removed at
length forty-nine.  Afterwards $\mathsf a,\mathsf R$
have a common survivor flag only at residue one.

Let $C_t$ denote the common survivor flag of
$\mathsf c_t$ and $\mathsf M_{t+1}$, paired at
length thirty-three.  Then
\begin{equation}\label{g:Q-C}
 Q_0=1\quad\Longrightarrow\quad C_1=0.
\end{equation}
\end{lemma}

\begin{proof}
Suppose a copy of $\mathsf a_2$ survives forty-nine.
The only possible subsequent incoming through
ninety-seven is the length-eighty-five arrow from
$\mathsf q$, with shift $(3,8)$.  But $\mathsf q$
remains only at residue zero, so it can hit only
$\mathsf a_0$.  The potential length-ninety-seven
outgoing has target $\mathsf v$, now absent.
Thus $\mathsf a_2$ is a cycle on page $97$.
The Euler copy $p^E$ is also still present: the
length-fifty-seven arrow is zero, and the possible
length-ninety-three incoming from $\mathsf y_0$
has target residue two, not zero.  The product
$p^E\mathsf a_2$ is consequently a required
$\mathsf f_2$ that is a product of two cycles at
ninety-seven, a contradiction.

Suppose instead that $\mathsf a_0$ survives forty-nine.
The detector $p^\dagger$ is present on page $85$;
its only earlier possible incoming was the zero
length-fifty-seven map.  If the hypothesized
$\mathsf a_0$ were hit at eighty-five, multiplying
by this detector would make its nonzero
$\mathsf f_2$ product an earlier boundary, contrary
to the requirement that every $\mathsf f_2$ remain
until ninety-seven.  Hence $\mathsf a_0$ remains
to ninety-three.  Nor can $p^\dagger$ be hit there
under this hypothesis: if $d_{93}z=p^\dagger$, then
\[
 d_{93}(z\mathsf a_0)=p^\dagger\mathsf a_0
\]
would again make a required $\mathsf f_2$ an
earlier boundary.  Both factors therefore reach
ninety-seven as cycles, and their product gives the
same contradiction.  Thus $\mathsf a_0$ and
$\mathsf a_2$ cannot survive forty-nine.  Their
one-dimensional targets $\mathsf R_0,\mathsf R_2$
are removed with them.  The common remaining flag
at residue one follows from the same pairing and
the unit symmetries through fifty-three.

Finally suppose $Q_0=1$.  A remaining
$\mathsf q_0$ stays to page $97$: its potential
eighty-five target $\mathsf a_0$ is absent.
If $C_1=1$, an incoming on $\mathsf c_1$ after
thirty-three and before ninety-seven, multiplied by $\mathsf q_0$, would
make the nonzero product $\mathsf q_0\mathsf c_1$
a required $\mathsf f_2$ that is an earlier
boundary.  Thus $\mathsf c_1$ must itself remain
to ninety-seven; then their product is a product
of cycles there, also impossible.  This proves
\eqref{g:Q-C}.
\end{proof}

\subsection{The last unit and the complete remaining arrow list}

After removing $\mathsf r,\mathsf v,\mathsf b,\mathsf d$,
the following table lists all candidate arrows modulo
$216$, before testing the $t$ residues of a proposed
partner:
\begin{equation}\label{g:remaining-mod216}
\begin{array}{c|l}
13&\mathsf q\longrightarrow\mathsf c\\
17&\mathsf A\longrightarrow\mathsf y,\quad
   \mathsf X\longrightarrow\mathsf M\\
29&\mathsf q\longrightarrow\mathsf X\\
33&\mathsf c\longrightarrow\mathsf M\\
49&\mathsf a\longrightarrow\mathsf R\\
53&\mathsf A\longrightarrow\mathsf x\\
57&\mathsf x\longrightarrow\mathsf p\\
61&\mathsf R\longrightarrow\mathsf f\\
85&\mathsf p\longrightarrow\mathsf c,\quad
   \mathsf q\longrightarrow\mathsf a\\
93&\mathsf y\longrightarrow\mathsf p\\
97&\mathsf f\longrightarrow\mathsf A\\
101&\mathsf p\longrightarrow\mathsf X\\
105&\mathsf X\longrightarrow\mathsf R\\
121&\mathsf c\longrightarrow\mathsf R\\
129&\mathsf x\longrightarrow\mathsf q\\
149&\mathsf M\longrightarrow\mathsf f\\
157&\mathsf p\longrightarrow\mathsf a\\
165&\mathsf y\longrightarrow\mathsf q\\
177&\mathsf a\longrightarrow\mathsf M.
\end{array}
\end{equation}
The table follows from the three explicit congruences
\eqref{g:l-congruence}--\eqref{g:four-congruence}.
For subsequent pages one uses the positive representative
of the length after the page in question, adding $216$
and $(4,24)$ as specified in
\eqref{g:copy-shifts}.  Thus the table includes all
later repetitions within the bound, rather than only
the displayed small lengths.

\begin{proposition}\label{g:unit313}
All translations in $\mathcal L$ are unit cycles through
page $E_{313}$.  The unit is hit at length $313$ from
the $\mathsf f$ translate
\begin{equation}\label{g:last-f}
 (-5,-36)+\mathcal L.
\end{equation}
There are no other arrows on the remaining labels at
that length, and the next page is zero.
\end{proposition}

\begin{proof}
We first make the units giving the translations explicit.
On $E_{53}$ put
\[
 \Omega=(C^3)^{-1}d_{53}(C^3).
\]
Equation~\eqref{g:omega-form} says that the logarithmic
differentials of $C^3$ and $u'$ are $\Omega$ and
$\kappa\Omega$.  Their common target is a
residue-field line, so $3\Omega=0$.  Hence the units
\begin{equation}\label{g:lattice-generators}
 U_0=C^9,\qquad U_1=C^{-3\kappa}u'
\end{equation}
and their inverses are cycles.  Choose any integer
representative of $\kappa$; changing it by three
changes $U_1$ by a power of $U_0$.
Their shifts $(9,0)$ and $(-3\kappa,1)$ generate
$\mathcal L$.  Together with $\mathcal D$, they
realize every required translation by actual
multiplication on the page.

At the end of length ninety-seven, all transfer
kernels of $\mathsf A$ have disappeared.  Only its
restriction lifts at $\mathcal L$ remain.  These
are products of surviving units, and their inverses
are cycles by derivation.  They have no outgoing
between fifty-three and ninety-seven; the incoming
at ninety-seven lands in their transfer kernel,
not in their nonzero restriction quotient.
In particular, passing to this quotient preserves
the identities $U_iU_i^{-1}=1$: restriction sends
these units to units on the still nonzero $H$ page,
so $1$ is not a length-ninety-seven boundary.
No splitting of the earlier $\mathsf A$ extension
is needed.

For a remaining $\mathsf A$ the only possible
outgoing before $313$ are to $\mathsf d,\mathsf y,
\mathsf x$ at $229,233,269$, respectively.  The
first target has been removed.  The second has
$t\equiv1\pmod3$, whereas $\mathsf y$ occurs
only at $t\equiv0$.  The third has
$t\equiv2\pmod3$, whereas $\mathsf x$ occurs
only at $t\equiv1$.  Thus all those outgoing
maps are excluded.  Formula~\eqref{g:A-incoming}
has no remaining incoming before $313$ and no
other one through the bound $469$.  The bound in
Proposition~\ref{det:euler-bound} forces a nonzero
$\mathsf f\to\mathsf A$ at $313$ hitting the
unit.  Its shift, computed by
\eqref{g:copy-shifts}, is $(5,36)$, giving
\eqref{g:last-f}.

No other arrow in \eqref{g:remaining-mod216}
has length $313$.  A differential hitting one
contracts the differential graded algebra:
if $dz=1$ and $dy=0$, then $d(zy)=y$ up to the
usual sign.  Its cohomology is therefore zero.
This proves the proposition, including the actual
page range of the unit symmetries.
\end{proof}

Work from now on modulo $\mathcal L$ and the permanent
$\mathcal D$ translation.  Each fully present label
of a fixed residue $t\bmod3$ has three copies.
Apart from the distinguished $\mathsf f$ in
\eqref{g:last-f}, and the already specified
$\mathsf f_2$ pairings with transfer kernels,
every copy must be removed by a nonzero arrow
before $313$.  Each is one $k$.  A nonzero arrow removes both
endpoints, so neither can participate on a later page.
On the same page, two consecutive nonzero maps between
one-dimensional entries are excluded by $d_R^2=0$.
Consequently the remaining problem is an ordinary
matching of these copies, uniform under the now
established unit translations.

\begin{lemma}\label{g:p1-matching}
One has $C_0=1$.  If $C_1=1$, all
$\mathsf c_0,\mathsf M_2$ copies pair at length
$249$, and all $\mathsf p_1,\mathsf c_1$ copies
pair at length $301$.  If $C_1=0$, all
$\mathsf p_1,\mathsf c_0$ copies pair at length
$85$, and neither $\mathsf c_1$ nor
$\mathsf M_2$ is present.
\end{lemma}

\begin{proof}
There is no incoming to $\mathsf p_1$ in the
remaining array.  For explicit verification,
the possible incoming lengths from
\eqref{g:remaining-mod216} are $57,93,273,309$:
their sources have, respectively, residues
$\mathsf x_0,\mathsf y_2,\mathsf x_2,
\mathsf y_1$, all absent.  The first arrow is
also zero by transfer.  The removed label
$\mathsf d$ would have supplied incoming at
$97$ or $313$, and supplies none now.
Its possible outgoing partners through $313$
are
\[
 \begin{array}{c|c}
 85&\mathsf c_0\\
 101&\mathsf X_0\text{ (absent)}\\
 157&\mathsf a_2\text{ (absent)}\\
 301&\mathsf c_1.
 \end{array}
\]
The potential removed-label target at $229$ is
$\mathsf b$, also absent.  Thus every
$\mathsf p_1$ must pair with $\mathsf c_0$ at
$85$ or $\mathsf c_1$ at $301$.

Similarly the remaining possibilities for
$\mathsf M_2$ are an outgoing at $149$ to
$\mathsf f_2$, already consumed at $97$, and
incoming at $177,233,249$ from
$\mathsf a_2,\mathsf X_0,\mathsf c_0$,
respectively.  Only the last is present.  The
earlier candidate at $105$ came from the
removed $\mathsf b$ label.  This list is
exhaustive by the same congruences.  Thus
$\mathsf M_2$, if present, must be hit from
$\mathsf c_0$ at $249$.

If $C_1=1$, the length-thirty-three pairing
left both $\mathsf c_1$ and $\mathsf M_2$.
The unique partner of every $\mathsf M_2$
forces $C_0=1$ and reserves every
$\mathsf c_0$ copy for length $249$; the
shift $(5,28)$ is a bijection on the copies.
Therefore all $\mathsf p_1$ must use
$\mathsf c_1$ at $301$.
If $C_1=0$, neither $\mathsf c_1$ nor
$\mathsf M_2$ occurs, and every
$\mathsf p_1$ must use $\mathsf c_0$ at
$85$, again forcing $C_0=1$.  This proves
both alternatives.  In particular
$\mathsf M_1$ was present after length
thirty-three in either case.
\end{proof}

\subsection{Three rows and one unmatched source}

We now discard the pairs fixed in Lemma~\ref{g:p1-matching},
the $\mathsf A$ transfer kernels and their required
$\mathsf f_2$ partners, and the unit targets themselves.
Retain a mark on the distinguished $\mathsf f$ source
that leaves to the unit at $313$.  The remaining
matching has three rows, indexed by
\[
 \mathcal R=\{(t,u)\in\Z^2:3\mid t\}/\mathcal L.
\]
For a row $s\in\mathcal R$, place the remaining copies
at the offsets displayed in Table~\ref{g:row-offsets}.
This is a list of their occurrences immediately
before processing the candidate pairings of lengths
greater than fifty-three.  Thus a later pairing is
recorded as a pair of entries even when it has already
been logically forced.

\begin{table}[htbp]
\centering
\setlength{\tabcolsep}{5pt}
\begin{tabular}{c|rrrrrr}
\toprule
label&$\mathsf y$&$\mathsf p_2$&$\mathsf p_0$&
 $\mathsf M_1$&$\mathsf f_0$&$\mathsf f_1$\\\midrule
$(t,u)$&$(0,0)$&$(2,10)$&$(6,34)$&
 $(13,70)$&$(12,62)$&$(16,86)$\\\bottomrule
\end{tabular}

\medskip
\begin{tabular}{c|rrrrrr}
\toprule
label&$\mathsf q_0$&$\mathsf X_1$&$\mathsf c_2$&
 $\mathsf M_0$&$\mathsf a_1$&$\mathsf R_1$\\\midrule
$(t,u)$&$(3,18)$&$(4,22)$&$(8,42)$&
 $(9,46)$&$(10,50)$&$(10,58)$\\\bottomrule
\end{tabular}
\caption{Offsets from the row index $s$.  The first six
entries occur in every row.  The last six occur in
pairs, with flags $Q_0,C_2$, and the survivor flag of
$\mathsf a_1$, respectively.}\label{g:row-offsets}
\end{table}

There is also the singleton $\mathsf x$ translate class.
Assign it offset $(1,4)$.  Its actual support is
$(1,6)+\mathcal L$ by \eqref{g:A-x-support}, so it
belongs to row
\begin{equation}\label{g:x-row}
 s_x=(0,2)+\mathcal L.
\end{equation}
The complete arrow list between the entries of
Table~\ref{g:row-offsets}, and $\mathsf x$, is
\begin{equation}\label{g:row-arrows}
\begin{array}{c|l}
61&\mathsf R_1\longrightarrow\mathsf f_0\\
85&\mathsf p_0\longrightarrow\mathsf c_2\\
93&\mathsf y\longrightarrow\mathsf p_2\\
101&\mathsf p_2\longrightarrow\mathsf X_1\\
121&\mathsf c_2\longrightarrow\mathsf R_1\\
129&\mathsf x\longrightarrow\mathsf q_0\\
149&\mathsf M_0\longrightarrow\mathsf f_0,\quad
     \mathsf M_1\longrightarrow\mathsf f_1\\
157&\mathsf p_0\longrightarrow\mathsf a_1\\
165&\mathsf y\longrightarrow\mathsf q_0\\
177&\mathsf a_1\longrightarrow\mathsf M_1\\
229&\mathsf q_0\longrightarrow\mathsf c_2\\
233&\mathsf X_1\longrightarrow\mathsf M_0\\
249&\mathsf c_2\longrightarrow\mathsf M_1\\
273&\mathsf x\longrightarrow\mathsf p_0\\
277&\mathsf R_1\longrightarrow\mathsf f_1\\
301&\mathsf p_2\longrightarrow\mathsf c_2,\quad
     \mathsf q_0\longrightarrow\mathsf a_1\\
309&\mathsf y\longrightarrow\mathsf p_0.
\end{array}
\end{equation}
The length-fifty-seven map has already been proved zero.
No other edge is suppressed in \eqref{g:row-arrows}:
it is obtained from \eqref{g:remaining-mod216} by the
residue conditions, removal of the pairs in
Lemma~\ref{g:p1-matching}, and the established absent
flags.  In particular all repetitions with length
less than $313$ are included.

Every displayed arrow preserves the row.  This can
be checked directly by subtracting the source
offset in Table~\ref{g:row-offsets} from the target
offset and applying \eqref{g:copy-shifts}.  To make
the entire finite calculation explicit, the shifts
are
\begin{equation}\label{g:row-shift-table}
\begin{array}{c|c@{\qquad}c|c}
R&(\Delta t,\Delta u)&R&(\Delta t,\Delta u)\\\hline
61&(2,4)&85&(2,8)\\
93&(2,10)&101&(2,12)\\
121&(2,16)&129&(2,14)\\
149&(3,16)&157&(4,16)\\
165&(3,18)&177&(3,20)\\
229&(5,24)&233&(5,24)\\
249&(5,28)&273&(5,30)\\
277&(6,28)&301&(6,32)\text{ or }(7,32)\\
309&(6,34)&&
\end{array}
\end{equation}
At $301$ the shifts are $(6,32)$ for
$\mathsf p_2\to\mathsf c_2$ and $(7,32)$ for
$\mathsf q_0\to\mathsf a_1$; the two arrows at
$149$ have the same shift.  For example,
$\mathsf x\to\mathsf p_0$ has difference
$(6,34)-(1,4)=(5,30)$, and
$\mathsf R_1\to\mathsf f_1$ has difference
$(16,86)-(10,58)=(6,28)$.  Thus there is no
possible differential between two distinct rows.

\begin{proposition}\label{g:kappa}
The coefficient in \eqref{g:lattice} is
$\kappa=-1\in\F_3$.  Equivalently,
\[
 \mathcal L=\{(t,u):t-3u\equiv0\pmod9\}.
\]
\end{proposition}

\begin{proof}
Each row of Table~\ref{g:row-offsets} begins with
six entries and then three optional pairs, so it
has even cardinality.  The row $s_x$ has one
additional entry, the singleton $\mathsf x$.
By Proposition~\ref{g:unit313}, every entry pairs
within its row except the marked $\mathsf f$
that leaves to $\mathsf A$ at $313$.  Removing
two endpoints preserves parity.  Therefore the
marked $\mathsf f$ must belong to the same row
as $\mathsf x$.

The marked source is $(-5,-36)+\mathcal L$ by
\eqref{g:last-f}.  It is an $\mathsf f_1$,
whose offset is $(16,86)$, so its row is
\[
 s_f=(-21,-122)+\mathcal L.
\]
Equality with \eqref{g:x-row} says
$(-21,-124)\in\mathcal L$, or
\[
 -21-372\kappa\equiv0\pmod9.
\]
Since $\kappa\in\F_3$, this is equivalent to
$\kappa=2=-1$.  The conclusion follows.
\end{proof}

\subsection{The known survivor and every higher detector}

\begin{proof}[Proof of Theorem~\ref{g:tail-exclusion}]
By Proposition~\ref{g:kappa}, copies are distinguished
by the key $t-3u$ modulo nine.  The three rows have
keys $0,3,6$.  The exceptional $\mathsf x$ row
$s_x=(0,2)$ has key $3$.

The detector $p^\dagger$ for $b_2$ has coordinates
$(-4,12)$ and hence key $5$.  Its $\mathsf p_2$
offset $(2,10)$ in Table~\ref{g:row-offsets} has
key $8$, so this detector lies in row $5-8=6$.
In that row the only available incoming on
$\mathsf p_2$ is the length-ninety-three arrow
from $\mathsf y$.  The $\mathsf x$ arrows
belong to row $3$, and all other incoming
labels have been removed.  Since $b_2$ is
known to survive, the detector cannot leave
by an outgoing differential while present.
It must nevertheless disappear before $313$
by Proposition~\ref{g:unit313}.  Consequently
the $\mathsf y$ of row $6$ hits this
$\mathsf p_2$ at length $93$.

Consider now $\mathsf p_0$ in that same row.
Its potential incoming from $\mathsf x$ at
$273$ belongs to row $3$ and cannot hit it.
Its only other incoming is the length-$309$
map from the row's $\mathsf y$, already
consumed at $93$.  Thus this $\mathsf p_0$
has no incoming.  It must disappear by an
outgoing differential; the full table permits
its outgoing at $85$ to $\mathsf c_2$ or
at $157$ to $\mathsf a_1$.  The argument does
not require a choice between these two
lengths. Figure~\ref{g:forcing-figure} summarizes this constraint.

\begin{figure}[htbp]
\centering
\begin{tikzpicture}[
  >=Stealth,
  every node/.style={font=\small},
  entry/.style={inner sep=5pt},
  arrowlabel/.style={fill=white,inner sep=2pt,align=center}
]
  \draw[rounded corners,black!35] (-0.55,-1.95) rectangle (10.25,1.25);
  \node[anchor=west,fill=white,inner sep=3pt] at (-0.3,1.25)
    {Row $6$};
  \node[entry] (y) at (0.4,0) {$\mathsf y$};
  \node[entry,align=center] (p2) at (4.2,0.65)
    {$\mathsf p_2$\\[-1pt]\scriptsize detector for $b_2$};
  \node[entry,align=center] (p0) at (4.2,-0.85)
    {$\mathsf p_0$\\[-1pt]\scriptsize detectors for $b_{\jmath}$, $\jmath\ge4$};
  \node[draw,rounded corners,align=center,inner sep=6pt] (out) at (8.8,-0.85)
    {one of\\[2pt]
     $d_{85}:\ \mathsf c_2$\\[-1pt]
     $d_{157}:\ \mathsf a_1$};
  \draw[->,thick] (y) -- node[arrowlabel,above,sloped]
    {$d_{93}$ forced} (p2);
  \draw[->,dashed] (y) -- (p0);
  \node[arrowlabel] at (1.5,-1.25) {$d_{309}$ unavailable};
  \draw[->,thick] (p0) -- node[arrowlabel,above]
    {outgoing forced} (out);
\end{tikzpicture}
\caption{The final constraint in row $6$.  The known survivor $b_2$
forces the $d_{93}$ pairing, so its source $\mathsf y$ cannot also
hit $\mathsf p_0$ at length $309$.  The possible $\mathsf x$ incoming
at length $273$ lies in row $3$.  Thus $\mathsf p_0$ must support
one of the two indicated outgoing differentials; the argument does
not choose between them.  The diagram records logical constraints,
not chronological order: length $85$ precedes length $93$.}
\label{g:forcing-figure}
\end{figure}


The $\mathsf p_0$ offset $(6,34)$ has key
$6-3\cdot34=3$ modulo nine.  In row $6$ it
therefore has total key zero.  Every
detector with $\jmath\ge4$ in
\eqref{g:detector-coordinates} has
\[
 t\equiv0\pmod3,\qquad
 t-3u=-40\cdot3^{\jmath-2}\equiv0\pmod9.
\]
It is exactly a $\mathsf p_0$ copy of this
key.  The unit translations of
Proposition~\ref{g:unit313} identify all such
copies through the relevant pages.  In
particular none can be an earlier incoming
Tate boundary, and each must have the forced
outgoing just proved.  This contradicts the
cycle property of the image of any surviving
sphere class.  Hence no $b_{\jmath}$ with
$\jmath\ge4$ survives.

For comparison, the $\jmath=3$ detector has
key $6$, whereas the $\jmath=2$ detector has
key $5$.  Thus the tail exclusion does not
decide $b_3$ by the same argument.  Its
survival and order-three representative
will be constructed separately in
Section~\ref{sec:cube}.
\end{proof}
